\chapter*{怎样读这一卷}
\addcontentsline{toc}{chapter}{怎样读这一卷}

本卷以九六〇年至一三六八年为主，同时写辽、党项与西夏、金、蒙古和元的形成、转换与区域治理；辽亡后的西辽、金蒙转换、西夏亡和元初征服也不只是附录。宋、辽、西夏、金、元及其竞争政权都是行动者，不能把宋朝年号当作唯一钟表。

请先对照各政权的原始年号、干支、月日和换算，沿着并行的年序追踪外交、战争、财政和迁徙；再阅读专题如何把交通、贸易、征发、家庭与制度接到这些事件上。因果阅读要分开长期条件、临近触发、可动员的资源和组织，以及结果如何反过来改变下一段时间线，而不以“正统”“夷夏”或亡国标签替代解释。

五史、《续资治通鉴长编》《建炎以来系年要录》《宋会要辑稿》及宋元政书提供不同的编年和制度入口；契约、文书、碑刻、考古以及契丹、女真、蒙古、藏、波斯、高丽和海域材料带来其他观察位置。诏令、统计、墓志或一批出土材料只能在其机构、地点、语言和保存条件所允许的范围内说话。

\texttt{event}、\texttt{turningpoint} 和 \texttt{causalchain} 分别帮助定位事件、转折和因果层次；\texttt{textwindow} 限定短引文的语境；\texttt{sourcenote}、\texttt{debate} 让不同材料与解释并列。使用它们回到前后的叙事，比较各政权怎样面对同一条道路、边界或市场，而非寻找单一中心的答案。
